本实验在本地编写攻击程序，并通过访问实验指导中给出的解密服务器完成 Padding Oracle 攻击。实验过程中，攻击者不需要知道 AES 密钥，只需要根据服务器返回的填充状态判断密文解密后的 PKCS\#7 填充是否合法。

\begin{table}[H]
    \centering
    \caption{实验环境}
    \begin{tabular}{|c|c|}
        \hline
        项目 & 配置或说明 \\
        \hline
        操作系统 & Windows 11 \\
        \hline
        编程语言 & Python 3 \\
        \hline
        开发工具 & Visual Studio Code \\
        \hline
        加密算法 & AES-128-CBC \\
        \hline
        填充标准 & PKCS\#7 \\
        \hline
        目标密文编号 & 5 号密文 \\
        \hline
        密钥组 & Key\_2 \\
        \hline
        解密接口 & \verb|http://10.102.33.67:8208/dec_2?data=| \\
        \hline
    \end{tabular}
\end{table}

其中，实验指导给出的解密服务器会对输入的密文进行 AES-128-CBC 解密，并检查解密后明文是否满足 PKCS\#7 填充规则。若填充正确，服务器返回 HTTP 200；若填充错误，服务器返回 HTTP 500 server error。因此，本实验程序通过构造不同的密文并观察服务器返回状态，逐步恢复目标密文对应的明文。

由于本小组抽到的是 5 号密文，而实验材料中 5 号密文属于 Key\_2 加密得到的密文，因此程序中访问的解密接口为 \verb|dec_2|。实验运行时需要保证本机能够访问学校实验服务器，通常需要处于校园网环境或连接学校 VPN。